\documentclass[a4paper, 10pt]{article}
\usepackage[margin=1in]{geometry}

\usepackage[T1]{fontenc}
\renewcommand*\familydefault{\sfdefault}
\usepackage{sfmath}
\usepackage{amsmath}





\begin{document}

\noindent \textbf{Individual Interim Report}\\
\noindent Tahmid Azam (ta549@cam.ac.uk)\\
\noindent Part IIA Project GG4: \textit{Neural Control with Adaptive State Estimation}\\
\noindent 22 May 2026

\vspace{1em}
\hrule
\vspace{1em}

% definining scope, what is estimation? How do we measure a good estimator?

\section*{Defining the parameter space}

To evaluate the estimator, we want to test its performance across a parameter space. Each parameter in should be: interpretable, communicating a particular characteristic of the behaviour of the LGSSM. A condition is defined by a set of parameters: the matrices defining the LGSSM ($\underline{A}$, $\underline{B}$, $\underline{C}$, $\underline{Q}$, and $\underline{R}$); the duration of the simulation, $T$; the number of trials, $n_\text{trials}$; the input, $u_t$; and the initial state, $x_0$.

We can derive additional parameters from these fundamental parameters. From the shapes of the matrices defining the LGSSM we find: the dimension of the latent states, $n_x$, from the shape of $\underline{A}$, $(n_x, n_x)$; the dimension of the inputs, $n_u$, from the shape of $\underline{B}$, $(n_x, n_u)$; and the dimension of the observations, $n_y$, from the shape of $\underline{C}$, $(n_y, n_x)$. These dimensions form new scalar parameters that we can vary.

We can also derive the spectral radius of $\underline{A}$, $\rho(\underline{A})$, which communicates the stability, stationarity, and memory of the LGSSM.

We can derive the controllability and observability Gramians, $\underline{W}_c$ and $\underline{W}_o$ respectively. We can use the determinant, trace, minimum eigenvalue, and the condition number.

We can derive the input-to-state gain

We can derive the observation gain

We can derive the observation ratio

We can directly vary the signal types and initial conditions.

Analytically generating a set of conditions that vary themselves uniformly across a high dimensional parameter space of these derived quantities would require a large amount of code; instead we propose a strategy to randomise the generated conditions, and keep randomly generating new conditions until the parameter space is effectively traversed. These conditions can then be saved to disk and reused

% Transforming the parameter space from matrices to more interpretable parameters

% How do we sample from this parameter space?
% Baseline then indepedently vary one axis
% Test every case, but this explodes, so use LHS?

% defining an evaluation process
% How can we vary signals? Signal generator module
% How can we vary noise? vary SNR on a continuum?
% How can we vary dimensions (high neurone, long duration, lots of timesteps)
% How can we vary the different matrices?

\end{document}